\section{Теоретическая часть}
\subsection*{Введение}

Отражательный клистрон
\footnote{Название <<клистрон>> происходит от греческого слова «клизо», что означает береговой волнорез.}, 
изучению которого посвящена лабораторная работа, предназначен для генерации электромагнитных колебаний СВЧ диапазона. Исторически именно генераторы клистронного типа позволили освоить диапазоны сантиметровых и дециметровых волн, в которых традицион­ные электронные лампы оказались неэффективными [1].

Генерация в клистронах осуществляется за счет преобразования кинетической энергии пучка электронов, ускоренных статическим электрическим полем, в энергию сверхвысокочастотных колебаний. Для получения СВЧ колебаний пучок электронов модулируется по плотности и проводится че­ рез резонатор клистрона, в котором переменный конвекционный ток пучка возбуждает колебания.

Неоднородный электронный поток в приборах клистронного типа созда­ется посредством так называемого динамического способа управления плот­ностью зарядов. При этом способе управления однородный электронный по­ток пропускается через переменное электрическое поле управляющего 
(груп­пирующего) устройства. Переменное поле воздействует на скорости электро­
нов в потоке, периодически ускоряя и замедляя их движение в зависимо­
сти от фазы высокочастотного поля, существующей в момент пролета элек­тронов через управляющее устройство. Возникающее различие в скоростях
электронов приводит к их группировке при последующем движении за пре­
делами управляющего устройства. При этом в определенной точке простран­
ства в определенный момент времени происходит образование электронного
сгустка. Поток приобретает пульсирующий характер, причем электронные
уплотнения и разрежения в данной точке пространства возникают с перио­дичностью, соответствующей частоте управляющего поля
\footnote{При использовании приемов т. н. электростатического управления основной результат управ­
	ления состоит в мо дуляции электронного потока по плотности непосредственно в области управ­
	ляющего поля [2, 3].}.

Характер процесса группировки электронов, координата и время образования сгустка в клистронах определяются условиями движения вне управляющего промежутка.

В пролетных клистронах группировка происходит при движении электронов по инерции в пространстве, свободном от внешних постоянных или переменных полей; группировка становится возможной за счет того, что «быстрые» электроны догоняют «медленные», вылетевшие из управляюще­го промежутка раньше «быстрых».

В отражательных клистронах электроны, вышедшие из группирующего устройства, движутся в постоянном тормозящем электрическом поле. Значение этого поля таково, что все электроны, вылетевшие из управляющего устройства, возвращаются назад. Группировка в отражательных клистронах происходит благодаря тому, что «быстрые» электроны находятся в пространстве группировки дольше «медленных», и на обратном пути оказывается возможной встреча «быстрых» электронов с медленными, вышедшими из управляющего промежутка позже «быстрых».

Принципиальная схема отражательного клистрона представлена на рис. \ref{fig:1}. Функции управляющего устройства и устройства, накапливающего энергию электромагнитных колебаний в отражательном клистроне, объединены здесь в тороидальном резонаторе.

\begin{figure}[h!]
	\centering
	\includegraphics[width=0.4\textwidth]{fig1}
	\caption{Идеализированная принципиальная схема клистрона: 1 — катод, 2 — резонатор, 3 — отражатель, 4 — вывод энергии, 5 — электронный поток, 6 — управляющий электрод}
	\label{fig:1}
\end{figure}

Электронный поток, эмитируемый катодом, ускоряется в промежутке между катодом и резонатором, после чего первый раз попадает в зазор резонатора, образованный двумя прозрачными для потока металлическими сетками. В зазоре происходит модуляция электронного потока по скорости. После выхода из резонатора электроны движутся в тормозящем поле отражающего электрода, возвращаются назад и повторно проходят зазор. При соответствующем выборе потенциалов на электродах клистрона сгусток электронов формируется в сечении зазора резонатора в момент времени, когда высокочастотное поле в зазоре тормозит возвращающиеся электроны, при этом происходит преобразование кинетической энергии электронов в энергию колебаний резонатора. Существование стационарных колебаний в клистроне оказывается возможным при компенсации потерь и резонаторе и нагрузке притоком энергии от модулированного электронного потока.

\subsection{Резонатор клистрона}

Резонатор в клистроне выполняет две функции: служит для модуляции электронного потока по скорости и для преобразования кинетической энергии пучка электронов в энергию электромагнитного ноля. Использовать для выполнении этих функций резонатор, геометрические размеры которого сравнимы с длиной волны основного колебания, возбуждаемого в них, невозможно из-за низкой эффективности взаимодействия электромагнитного поля и пучка электронов, пролетающих через этот резонатор
\footnote{Следует отметить, что существует способ группировки электронного потока по плотности, при котором модулируемый поток проводится через область, в которой фаза управляющего поля успевает многократно измениться за время пролета электрона. При использовании этого способа скорости электронов на выходе из управляющего поля окатываются практически одинаковыми, а электронный поток промодулированным по плотности [3].}. Действительно. если средняя скорость электронов $v$ значительно меньше скорости света $с$, то время их пролета $t = a/v$ через резонатор, имеющий форму ку­ ба со стороной а, намного превосходит период колебания поля $T =\sqrt{2}a/c$. Электрон при движении через такой резонатор то ускоряется, то тормозится, поэтому обмен энергией между ним и полем незначителен. Время пролета можно уменьшить, взяв вместо куба призму с достаточно малой высотой:
при неизменной низшей частоте колебаний время пролета электрона через резонатор уменьшается. Однако сокращение объема резонатора приводит к
уменьшению его добротности. Использование в генераторе низкодобротной колебательной системы не позволяет производить эффективное преобразование кинетической энергии электронов в энергию электромагнитных колебаний и не обеспечивает стабильности частоты генератора.

Для увеличения объема резонатора и его добротности необходимо, оставляя длину пролетного промежутка малой, создавать дополнительные резервуары энергии. Именно такой резонатор используется в отражательном клистроне.

Схема резонатора представлена на рис. \ref{fig:2}. Геометрические размеры резонатора удовлетворяют неравенствам $a,b \ll \lambda /4,l \ll \lambda /4,d \ll l$, где $\lambda$ — длина волны, соответствующая частоте основной моды резонатора. По­скольку размеры резонатора много меньше длины волны на частоте основ­
ного колебания, резонатор можно рассматривать как колебательный контур с сосредоточенными параметрами.

При этом в резонаторе можно выделить емкостной объем, заключенный между сетками резонатора, и индуктивный объем, т. е. собственно тороидальную часть.

\begin{figure}[H]
	\centering
	\includegraphics[width=0.4\textwidth]{fig2}
	\caption{Поперечное сечение тороидального резонатора: 1 — емкостной объем резонатора, 2 — индуктивный объем резонатора}
	\label{fig:2}
\end{figure}

При наличии потерь энергии поля в резонаторе, возникающих из-за конечной проводимости стенок и отвода мощности в нагрузку, свободные колебания резонатора являются затухающими, при этом в рамках метода комплексных амплитуд удобно вводить комплексную частоту $\omega = w^{\prime}+i\omega ^{\prime \prime}$. Ве­
личина $\omega ^{ \prime \prime }$ может быть выражена через добротность резонатора, которая, как правило, измеряется экспериментально.

В клистронах применяются резонаторы, допускающие механическую перестройку резонансной частоты. Наиболее распространены два способа перестройки: индуктивный и емкостной. Индуктивная перестройка осуществляется введением в область магнитного поля металлических поршней, кото­рые, уменьшая занятый полем объем, уменьшают индуктивность резонатора и тем самым повышают его резонансную частоту. Емкостная перестройка может быть осуществлена путем деформации гибкой мембраны. Сближение сеток резонатора, обеспечиваемое прогибом мембраны, ведет к увеличению емкости зазора и уменьшению резонансной частоты. Механические способы перестройки обеспечивают отстройку от основной частоты на 20-25\%.

\subsection{Модуляция скорости электронов в пучке}

Рассмотрим процесс модуляции электронного потока по скорости. В установившемся режиме колебания в резонаторе близки к гармоническим, поле в зазоре имеет вид 
$E = E _ { z } = E _ { m } \text { sin } \omega t$, 
где $E _ { m } = \frac { U _ { m } } { d }$, $U _ { m }$ амплитуда напряжения на зазоре,$E _ { z }$ проекция электрического ноля на ось $z$,
$d$ - ширина зазора, $\omega$ частота стационарных колебаний. Найдем приращение энергии электрона, прошедшего через зазор. Кинетическая энергия,
приобретаемая одиночным электроном при прохождении пути $dz$ внутри за­зора, равна работе силы электрического поля: 
$d W = - e \frac { U _ { m } } { d } \sin \omega t \, dz$, 
где $e$ - модуль заряда электрона (знак заряда учтен в явном виде)
\footnote{Если напряжение (потенциал второй сетки относительно первой) $U = -U_m \sin \omega t$ положительно, оно ускоряет электрон, движущийся в положительном направлении оси z.}. 
Электрон, прошедший расстояние между катодом и резонатором, имеет кинетическую энергию 
$\frac { m v _ { 0 } ^ { 2 } } { 2 } = e U _ { 0 }$ , где $U_0 > 0$-- потенциал резонатора относительно ка­тода, $v_0$ — скорость электрона, $m$ — его масса. Отсюда при условии $v _ { 0 } \ll c$ имеем
\footnote{Пренебрежение релятивистскими поправками возможно до значений $U_0$ порядка нескольких
	десятков киловольт.} 
$v _ { 0 } = \sqrt { \frac { 2 e U _ { 0 } } { m } }$

При выполнении условия 
$\frac { U _ { m } } { U _ { 0 } } \ll 1$ возмущения скорости $v_0$ под воздействием высокочастотного поля незначительны. Поэтому если $t_0$ — время про­
хождения электроном центра зазора, то время его нахождения в точке с координатой $z$ есть $t = t _ { 0 } + \frac { z } { v _ { 0 } }$. Полное приращение энергии имеет вид

\begin{equation}
	{ \Delta W = \int\limits _ { - d / 2 }^{{ + d / 2 }} \frac { - e U _ { m } } { d } \sin \left( \omega t _ { 0 } + \frac { \omega z } { v _ { 0 } } \right) d z = } \\ 
	{ - e U _ { m } \sin \omega t _ { 0 } \frac { \sin \left( \theta _ { 3 } / 2 \right) } { \theta_ { 3 } / 2 } = - e M U _ { m } \sin \omega t _ { 0 }, } 
\end{equation}

где $\theta _ { 3 } = \frac { \omega d } { t _ { 0 } }$ - невозмущенный угол пролета электрона через модулирующий зазор, $M = \frac { \sin \left( \theta _ { 3 } / 2 \right) } { \theta _ { 3 } / 2 }$ - коэффициент взаимодействия электронного потока с полем зазора. Полная кинетическая энергия электрона, вошедшего в зазор с начальной скоростью на выходе из него может быть представлена в виде:
\begin{equation}
	W = \frac { m v ^ { 2 } } { 2 } = \frac { m v_0 ^ { 2 } } { 2 } + \Delta W = e U _ { 0 } + \Delta W.
\end{equation}
Таким образом, скорость электрона на выходе из зазора оказывается равной
\begin{equation}
	v = \sqrt { \frac { 2 e } { m } U _ { 0 } } \left( 1 - \frac { M U _ { m } } { U _ { 0 } } \sin \omega t _ { 0 } \right) ^ { 1 / 2 }.
\end{equation}
При выполнении условия $\frac { U _ { m } } { U _ { 0 } } \ll 1$ можно разложить выражение для $v$ в ряд Тейлора и ограничиться двумя первыми членами ряда:
\begin{equation}
	v = \sqrt { \frac { 2 e } { m } U _ { 0 } } \left( 1 - \frac { 1 } { 2 } \frac { M U _ { m } } { U _ { 0 } } \sin \omega t _ { 0 } + \ldots \right) \approx v _ { 0 } - v _ { 1 } \sin \omega t _ { 0 },
\end{equation} где $v _ { 1 } = \frac { M U _ { m } } { 2 U _ { 0 } } v _ { 0 }$
Из полученного выражения видно, что скорость электро­на на выходе из зазора определяется фазой поля, существовавшей в момент прохождения им центра зазора. Наибольшая амплитуда скоростной модуляции $v_1$ достигается при стремлении коэффициента взаимодействия $M$ к единице, что выполняется при стремлении $\theta_3$ к нулю.

\subsection{Модуляция электронного пучка по плотности}

В пространстве между резонатором и отражателем электроны двигаются в статическом тормозящем поле. Уравнение движения имеет вид 
$mz''- e \frac { U _ { 0 } - U _ { \text{отр} } } { L }$ где $U_{\text{отр}} < 0$ — напряжение на отражателе, $L$ — расстояние между резонатором и отражателем. Интегрируя первый раз уравнение движения и учитывая выражение для скорости $v$ при выходе из зазора, имеем
\begin{equation}
	z ^ { \prime } = v - \frac { e } { m } \frac { U _ { 0 } - U _ { \text { отр } } } { L } \left( t - t ^ { \prime } \right)
\end{equation}
где $t'$ — момент выхода электрона из зазора. Через t в дальнейшем будем обозначать момент, когда тот же электрон возвращается в плоскость второй сетки. Повторное интегрирование уравнения (5) дает
\begin{equation}
	z = v \left( t - t ^ { \prime } \right) - \frac { e } { m } \frac { U _ { 0 } - U _ { \text { отр } } } { L } \frac { \left( t - t ^ { \prime } \right) ^ { 2 } } { 2 } + \frac { d } { 2 }
\end{equation}

Время пролета электрона в пространстве группировки можно найти из условия $z = d / 2$ при $t=t''$, откуда следует, что 
$t ^ { \prime \prime } - t ^ { \prime } = 0$ и 
$\frac { e } { m } \frac { U _ { 0 } - U _ { \text { отр } } } { L }\cdot \frac { \left( t ^ { \prime \prime } - t ^ { \prime } \right) } { 2 v } = 1$
Первое решение $(t' = t'' )$ соответствует моменту вылета электрона из зазора, второе дает время его пролета в тормозящем поле:
\begin{equation}
	t ^ { \prime \prime } - t ^ { \prime } = \frac { 2 m } { e } \frac { v L } { L _ { 0 } - U _ { \text { отр } } }
\end{equation}

При выполнении условия $U_{ m } / U _{ 0 } \ll 1$ время пролета в зазоре определяется скоростью $v_0$, поэтому связь времени вылета электрона из зазора со временем его нахождения в центре зазора можно приближенно записать в виде
$t'' = t _ { 0 } + d / \left( 2 v _ { 0 } \right)$. Таким образом, подставляя выражение для скорости (4) в (7), получим
\begin{equation}
	t ^ { \prime \prime } = t ^ { \prime } + \frac { 2 m L } { e \left( U _ { 0 } - U _ { \text { отр } } \right) } \left( v _ { 0 } - \frac { M U _ { m } } { 2 U _ { 0 } } v _ { 0 } \sin \left( \omega t ^ { \prime } - \frac { \omega d } { 2 v _ { 0 } } \right) \right)
\end{equation}

Из найденного выражения видно, что время пролета электронов в пространстве группировки зависит от фазы высокочастотного напряжения, существующего в момент пролета электроном середины зазора, и от амплитуды этого напряжения.

На рис.\ref{fig:3} представлены примеры пространственно-временных диаграмм(зависимостей координат электронов от времени) для таких значений потенциалов на резонаторе и отражателе, при которых электроны, вышедшие из резонатора в различные моменты времени, возвращаются в него одновремен­но, образуя сгусток. Сверхвысокочастотное поле в резонаторе в рассматриваемых в случаях в момент времени пролета сгустка максимально и является для него тормозящим. Уменьшение кинетической энергии электронов при этом приводит к возрастанию энергии СВЧ поля. Разумеется, при других значениях потенциалов на электродах клистрона сгусток может сформироваться и вне зазора резонатора, при этом эффективная передача энергии от пучка полю становится невозможной.

Домножим соотношение (8) на ы и введем следующие величины:
\begin{equation}
	\theta _ { \text{г} } = \frac { 2 m } { e } \frac { v _ { 0 } \omega L } { U _ { 0 } - U _ { \text{ oтр } } }
\end{equation} — угол пролета электрона в пространстве группировки,
\begin{equation}
	X = \theta _ { \text{г} } \frac{M U _ { m } } { 2 U _ { 0 } }
\end{equation} — параметр группировки. В результате получим выражение, связывающее время возвращения электрона в плоскость второй сетки $(t'' )$ со временем выхода электрона из зазора $(t')$:
\begin{equation}
	\omega t ^ { \prime \prime } = \omega t ^ { \prime } + \theta _ { r } - X \sin \left( \omega t ^ { \prime } - \frac { \theta _ { 3 } } { 2 } \right)
\end{equation}
Соотношение (11) является исходным для нахождения конвекционного тока пучка электронов.

\begin{figure}[h!]
	\centering
	\includegraphics[width=0.4\textwidth]{fig3}
	\caption{Пространственно-временные диаграммы движения электронов при двух значениях оптимального времени пролета $\tau$ в пространстве группировки ($z=0$ координата, соответствующая середине зазора)}
	\label{fig:3}
\end{figure}